\documentclass[11pt]{article}
\usepackage[margin=1in]{geometry}
\usepackage{amsmath}
\usepackage[hidelinks]{hyperref}
\setlength{\parindent}{0pt}
\setlength{\parskip}{4pt}

\begin{document}

\begin{center}
{\large\bfseries CEE292 Project Proposal: Improved PINNs for Option Pricing}\\[2pt]
Sondre Rogde
\end{center}

\textbf{Contribution:} We implement the paper Y. Bai, T. Chaolu, S. Bilige, ``The application of improved physics-informed neural network (IPINN) method in finance,'' \emph{Nonlinear Dynamics}, 2022 in code.
We implement both the standard PINN and their proposed improved PINN (IPINN).

\textbf{Data and benchmark:} No market data is needed. 
Training data consist of collocation, initial and boundary points sampled from the Black-Scholes and Ivancevic PDEs.
We evaluate on a test grid against the closed-form Black-Scholes price and the analytic Ivancevic solution.
We use relative $L^2$ error and maximum absolute error. 
We do this over 5 seeds and report mean $\pm$ standard deviation.

\end{document}
